\section{Decomposition}
\label{sec:theory}

Write $\atoms{\varphi}$ for the atomic propositions and data variables occurring
in $\varphi$, and $\clocks{\varphi}$ for the formula-clocks of the timed
modalities occurring in it. For a set $C$ of clocks, $\proj{Z}{C}$ is the zone
obtained by existentially quantifying away the clocks outside $C$.

\begin{definition}[Split]\label{def:split}
A node $N=\tup{\Gam}{\Ctx}{Z}{\obl}{\Dl}$ with pursuit set $P$ \emph{splits}
into $N_1,\dots,N_r$ ($r\ge 2$) when
\begin{enumerate}[label=\textup{(S\arabic*)},ref=\textup{(S\arabic*)}]
\item\label{S1} \Gam, \Ctx and $P$ are partitioned among the $N_i$, and
      $\atoms{N_i}\cap\atoms{N_j}=\emptyset$ for $i\ne j$;
\item\label{S2} every clock constrained in $Z$ belongs to exactly one
      $\clocks{N_i}$;
\end{enumerate}
and $N_i$ carries $\proj{Z}{\clocks{N_i}}$, the deadlines on its own clocks, and
the phase of $N$.
\end{definition}

There is deliberately no condition on the zone, and this is worth dwelling on
because the natural guess is that there should be. A reset leaves finite
difference constraints in the matrix, so $Z$ can relate $x$ and $y$ even when
the modalities owning them share no atom; it is tempting to require $Z$ to be
the product of its projections and to refuse the split otherwise. That condition
is sufficient and unnecessary --- Remark~\ref{rem:noproduct} --- and we
implemented it before noticing.

\begin{theorem}[Decomposition]\label{thm:decomp}
If $N$ splits into $N_1,\dots,N_r$, then $N$ is realisable if and only if every
$N_i$ is.
\end{theorem}

\begin{proof}
($\Rightarrow$) Let $\mathcal T$ realise $N$ at location $s$. Fix $i$ and let
$\mathcal T_i$ be $\mathcal T$ with its labelling restricted to $\atoms{N_i}$
and its clocks to $\clocks{N_i}$. Every formula of $N_i$ mentions only atoms of
$\atoms{N_i}$ and is evaluated against clocks of $\clocks{N_i}$, and the
satisfaction relation inspects a configuration only through those two, so the
truth of each is unchanged by the restriction; the runs and their delays are
unchanged, so \Runsinf is the same set and the path quantifiers range over it
unchanged. Each $\nu\in\proj{Z}{\clocks{N_i}}$ extends by definition of the
projection to some $\nu'\in Z$, and $\mathcal T,(s,\nu')$ satisfies the formulas
of $N_i$, hence so does $\mathcal T_i,(s,\nu)$.

($\Leftarrow$) Let $\mathcal T_i$ be induced by a timed automaton
$\mathcal A_i$ over $\clocks{N_i}$ realising $N_i$. The clock sets are pairwise
disjoint, so form the product $\mathcal A=\prod_i\mathcal A_i$: locations are
tuples, the invariant of a tuple is the conjunction of the components'
invariants, the labelling is the union of theirs --- consistent, the atom sets
being disjoint --- and edges interleave, each edge of one component lifting to
an edge of the product that leaves the other components' locations and clocks
untouched. Delays synchronise, and are admissible because each component's
invariant permits them.

\emph{Runs project and lift.} A run of $\mathcal A$ projects to a run of each
$\mathcal A_i$ with the same delays, so it is non-Zeno exactly when its
projections are. Conversely a run of one component and arbitrary runs of the
others merge into a run of the product projecting to each: refine the delay
sequences to their common refinement and interleave the discrete steps
consistently with time. Non-blockingness of the components gives
non-blockingness of the product, so \Runsinf is non-empty everywhere.

\emph{Truth is componentwise.} By induction on $\varphi$ with
$\atoms{\varphi}\subseteq\atoms{N_i}$. Literals are the definition of the
product labelling and the boolean cases are immediate. For a formula under \Ev,
take the witnessing run of $\mathcal A_i$ and lift it against arbitrary runs of
the others; for one under \Av, take any run of the product and project it, where
the hypothesis applies. A timed modality adds only a constraint on the elapsed
time of a position, which projection preserves because delays are shared.

\emph{The zone.} Let $\nu\in Z$. By \ref{S2} each coordinate belongs to exactly
one component, so $\nu$ decomposes as $(\nu_1,\dots,\nu_r)$ with
$\nu_i\in\proj{Z}{\clocks{N_i}}$, and $(s_i,\nu_i)$ realises $N_i$'s formulas.
By the previous paragraph the product configuration satisfies every formula of
every $N_i$, that is, all of $\Gam\cup\Ctx$. As $\nu$ was arbitrary,
$\mathcal A$ realises $N$. \qed
\end{proof}

\begin{remark}[No zone-product condition is needed]\label{rem:noproduct}
Let $Z$ be constrained by $x-y<1$ with $x$ and $y$ in different components. Both
projections are universal and $Z$ is not, so $Z$ is a proper subset of the
product of its projections. Splitting is still sound. Forwards, a valuation of a
projection extends to one of $Z$ and a component's formulas read only its own
clocks. Backwards, each component is asked to be realisable over its
\emph{projection}, which is the larger demand: $Z\subseteq\bigcap_i
\proj{Z}{\clocks{N_i}}$, so a product model meeting the projections meets $Z$ a
fortiori.

Dropping the condition costs nothing and buys nothing measurable: across the
whole suite it never once blocked a split that \ref{S1} and \ref{S2} allowed,
and node counts are identical with and without it. We report it because the
condition is the natural first guess and a reader will otherwise wonder why it
is absent.
\end{remark}

\begin{corollary}\label{cor:search}
If $N$ splits, its verdict in the tableau is the conjunction of the verdicts of
its components.
\end{corollary}

\begin{proof}
The calculus is sound and complete, so a node lies in the greatest successful
marking exactly when it is realisable; apply Theorem~\ref{thm:decomp}. \qed
\end{proof}
